\documentclass[12pt]{article}
\usepackage{amsmath,amssymb,amsthm}
\usepackage{booktabs,geometry,hyperref,cite}

\geometry{margin=1in}

\newtheorem{theorem}{Theorem}
\newtheorem{lemma}{Lemma}
\newtheorem{assumption}{Assumption}
\newtheorem{remark}{Remark}

\newcommand{\mc}{\mathcal}
\newcommand{\mbb}{\mathbb}
\newcommand{\E}{\mathbb{E}}
\newcommand{\Var}{\mathrm{Var}}
\newcommand{\Cov}{\mathrm{Cov}}
\newcommand{\tr}{\mathrm{tr}}
\newcommand{\ELPPD}{\mathrm{ELPPD}}

\begin{document}

\title{Double Asymptotic Variance and Bias of the $K$-fold ELPPD Estimator under High-Dimensional Bernstein-von Mises}
\author{Collaborative Research Draft}
\date{\today}
\maketitle

\begin{abstract}
We establish the double asymptotic theory ($S \to \infty$ followed by $n \to \infty$) of the self-normalized importance sampling (SNIS) estimator of the $K$-fold Expected Log Pointwise Predictive Density (ELPPD) computed from Markov chain Monte Carlo (MCMC) samples. Under a high-dimensional Bernstein-von Mises (BvM) framework and polynomial ergodicity conditions on the MCMC sampler, we derive the exact scaling of the asymptotic variance and the bias of its estimator. We prove that the first-order MCMC fluctuations vanish due to the MLE score equations, and the leading-order variance is determined by a quadratic form of the MCMC draws. The resulting asymptotic variance scales as $\frac{2 d(n)(1 - 2/K(n))}{(1 - 1/K(n))^2} \gamma_{\mathrm{mcmc}}$, where $d(n)$ is the parameter dimension and $\gamma_{\mathrm{mcmc}}$ is the MCMC efficiency factor.
\end{abstract}

\section{Introduction and Notation}
Let $y^{(n)} = (y_1, \dots, y_n)^\top \in \mathcal{Y}^n$ be a vector of $n$ independent observations drawn from the true data-generating distribution $P_{\theta_0}$ with density $p(y \mid \theta_0)$, where $\theta_0 \in \text{int}(\Theta)$ and $\Theta \subset \mathbb{R}^{d(n)}$ is the parameter space. The parameter dimension $d(n) = \dim(\theta)$ grows as a function of the sample size $n$. Let $\pi(\theta)$ be the prior distribution on $\Theta$.

Let $\kappa: \{1, \dots, n\} \to \{1, \dots, K(n)\}$ define the fold partition. The training set for fold $k = \kappa(i)$ is denoted by $\mathcal{D}_{-k}$, which has size $n - n_k = n(1 - 1/K(n))$, where $n_k \approx n/K(n)$ is the size of the omitted fold. The fold-specific posterior distribution is denoted by $p(\theta \mid y_{\mathcal{D}_{-\kappa(i)}})$.

Let $\{\theta^{(s)}\}_{s=1}^S$ be a Markov chain with stationary distribution $q(\theta)$, serving as the proposal distribution. For each observation $i \in \{1, \dots, n\}$, the unnormalized importance weights and weighted likelihoods are defined as:
\begin{equation}
w_{\kappa(i)}(\theta) = \frac{p(\theta \mid y_{\mathcal{D}_{-\kappa(i)}})}{q(\theta)}, \quad u_i(\theta) = p(y_i \mid \theta) w_{\kappa(i)}(\theta)
\end{equation}
The true fold-specific posterior predictive density is:
\begin{equation}
\mu_i := p(y_i \mid y_{\mathcal{D}_{-\kappa(i)}}) = \int p(y_i \mid \theta) p(\theta \mid y_{\mathcal{D}_{-\kappa(i)}}) d\theta = \frac{\mathbb{E}_q[u_i(\theta)]}{\mathbb{E}_q[w_{\kappa(i)}(\theta)]}
\end{equation}
and its self-normalized importance sampling (SNIS) estimator is:
\begin{equation}
\widehat{\mu}_i = \frac{\frac{1}{S}\sum_{s=1}^S u_i(\theta^{(s)})}{\frac{1}{S}\sum_{s=1}^S w_{\kappa(i)}(\theta^{(s)})}
\end{equation}
The estimator of the $K$-fold Expected Log Pointwise Predictive Density (ELPPD) is:
\begin{equation}
\widehat{\text{ELPPD}} = \sum_{i=1}^n \log \widehat{\mu}_i
\end{equation}
while the true mathematical $K$-fold ELPPD is $\text{ELPPD} = \sum_{i=1}^n \log \mu_i$.

\section{High-Dimensional BvM CLT for Fold-Specific Estimators}
We establish the joint asymptotic distribution of the fold-specific estimators as $n \to \infty$.

\begin{assumption}[Regularity of Likelihood and Prior]\label{ass:regularity}
We assume standard regularity conditions for growing parameter dimensions $d(n) \to \infty$ as $n \to \infty$ \cite{ghosal1999asymptotic, spokoiny2012parametric}:
\begin{enumerate}
    \item The dimension satisfies $d(n)^2 / n \to 0$ as $n \to \infty$.
    \item The log-likelihood function is twice continuously differentiable, and there exists a function $L(y)$ such that:
    \[
    \left\| \nabla^2_\theta \log p(y \mid \theta) - \nabla^2_\theta \log p(y \mid \theta_0) \right\|_{\text{op}} \leq L(y) \|\theta - \theta_0\|_2
    \]
    where $\mathbb{E}_{\theta_0}[L(Y)^2] \leq C_L < \infty$.
    \item The Fisher information matrix $I(\theta_0) = -\mathbb{E}_{\theta_0}[\nabla^2_\theta \log p(Y_1 \mid \theta_0)]$ is positive definite with eigenvalues bounded: $c_0 I_d \le I(\theta_0) \le C_0 I_d$.
\end{enumerate}
\end{assumption}

\begin{lemma}[Joint BvM CLT for Fold-Specific Estimators]\label{lem:joint_bvm}
Under Assumption~\ref{ass:regularity}, the fold-specific maximum likelihood estimators $\hat{\theta}_{-k}$ for $k \in \{1, \dots, K\}$ satisfy:
\[
\sqrt{n}\big(\hat{\theta}_{-k} - \theta_0\big) \xrightarrow{d} \mathcal{N}\left( 0, \frac{1}{1 - 1/K(n)} I(\theta_0)^{-1} \right) \quad \text{as } n \to \infty
\]
and the covariance between estimators for different folds $k \neq l$ is:
\begin{equation}
\text{Cov}\left( \hat{\theta}_{-k}, \hat{\theta}_{-l} \right) = \frac{n - n_k - n_l}{(n - n_k)(n - n_l)} I(\theta_0)^{-1} + o(1/n)
\end{equation}
which, for balanced fold sizes $n_k = n_l = n/K(n)$, is:
\begin{equation}
\text{Cov}\left( \hat{\theta}_{-k}, \hat{\theta}_{-l} \right) = \frac{1}{n} \frac{1 - 2/K(n)}{(1 - 1/K(n))^2} I(\theta_0)^{-1} + o(1/n)
\end{equation}
\end{lemma}

\begin{proof}
By standard asymptotic expansions for M-estimators in growing dimensions \cite{vdv1998asymptotic}, the training MLE $\hat{\theta}_{-k}$ satisfies:
\[
\hat{\theta}_{-k} - \theta_0 = I(\theta_0)^{-1} \frac{1}{n - n_k} \sum_{j \in \mathcal{D}_{-k}} \psi_j + o_p(1/\sqrt{n})
\]
where $\psi_j = \nabla_\theta \log p(y_j \mid \theta_0)$ are independent with $\mathbb{E}[\psi_j] = 0$ and $\text{Var}(\psi_j) = I(\theta_0)$.
The variance of the sum is:
\[
\text{Var}\left( \hat{\theta}_{-k} - \theta_0 \right) = I(\theta_0)^{-1} \left( \frac{1}{(n - n_k)^2} \sum_{j \in \mathcal{D}_{-k}} \text{Var}(\psi_j) \right) I(\theta_0)^{-1} = \frac{1}{n - n_k} I(\theta_0)^{-1} = \frac{1}{n(1 - 1/K)} I(\theta_0)^{-1}
\]
Since we scale the estimator difference by the factor $\sqrt{n}$, we multiply the variance by $n$:
\[
\text{Var}\left( \sqrt{n}\big(\hat{\theta}_{-k} - \theta_0\big) \right) = n \, \text{Var}\left( \hat{\theta}_{-k} - \theta_0 \right) = \frac{1}{1 - 1/K} I(\theta_0)^{-1}
\]
For two different folds $k \neq l$, the training sets are $\mathcal{D}_{-k}$ and $\mathcal{D}_{-l}$. The intersection of the training sets contains $n - n_k - n_l$ observations. Since folds $k$ and $l$ are disjoint, the covariance of their score sums is:
\[
\text{Cov}\left( \sum_{j \in \mathcal{D}_{-k}} \psi_j, \sum_{j \in \mathcal{D}_{-l}} \psi_j \right) = \sum_{j \in \mathcal{D}_{-k} \cap \mathcal{D}_{-l}} \text{Var}(\psi_j) = (n - n_k - n_l) I(\theta_0)
\]
Therefore, the covariance of the estimators is:
\[
\text{Cov}\left( \hat{\theta}_{-k}, \hat{\theta}_{-l} \right) = I(\theta_0)^{-1} \frac{n - n_k - n_l}{(n - n_k)(n - n_l)} I(\theta_0) I(\theta_0)^{-1} = \frac{n - n_k - n_l}{(n - n_k)(n - n_l)} I(\theta_0)^{-1}
\]
For equal fold sizes $n_k = n_l = n/K(n)$, this simplifies exactly to:
\[
\text{Cov}\left( \hat{\theta}_{-k}, \hat{\theta}_{-l} \right) = \frac{1}{n} \frac{1 - 2/K(n)}{(1 - 1/K(n))^2} I(\theta_0)^{-1}
\]
This completes the proof.
\end{proof}

\section{BvM Expansion of the Pointwise Predictive Density}
We derive the asymptotic form of the posterior predictive density $\mu_i$ under BvM.

\begin{theorem}[Pointwise Predictive Density Expansion]\label{thm:ppd_expansion}
Under Assumption~\ref{ass:regularity}, the logarithm of the fold-specific posterior predictive density satisfies:
\[
\log \mu_i = \log p(y_i \mid \hat{\theta}_{-\kappa(i)}) + \frac{1}{2 n (1 - 1/K)} \text{tr}\left[ \left( \nabla^2_\theta \log p(y_i \mid \theta_0) + \nabla_\theta \log p(y_i \mid \theta_0) \nabla_\theta \log p(y_i \mid \theta_0)^\top \right) I(\theta_0)^{-1} \right] + o_p\left( \frac{1}{n} \right)
\]
where $\hat{\theta}_{-\kappa(i)}$ is the MLE of the training set.
\end{theorem}

\begin{proof}
The predictive density is $\mu_i = \int p(y_i \mid \theta) p(\theta \mid y_{\mathcal{D}_{-\kappa(i)}}) d\theta$.
By the high-dimensional BvM theorem, we replace the training posterior $p(\theta \mid y_{\mathcal{D}_{-\kappa(i)}})$ with its Gaussian approximation $\mathcal{N}(\hat{\theta}_{-\kappa(i)}, \Sigma_{-\kappa(i)})$, where $\Sigma_{-\kappa(i)} = \frac{1}{n(1 - 1/K)} I(\theta_0)^{-1}$.

We Taylor-expand $p(y_i \mid \theta)$ around the training MLE $\hat{\theta}_{-\kappa(i)}$:
\[
p(y_i \mid \theta) = p(y_i \mid \hat{\theta}_{-\kappa(i)}) \exp\left( u^\top (\theta - \hat{\theta}_{-\kappa(i)}) + \frac{1}{2} (\theta - \hat{\theta}_{-\kappa(i)})^\top M (\theta - \hat{\theta}_{-\kappa(i)}) + o_p(\|\theta - \hat{\theta}_{-\kappa(i)}\|^2) \right)
\]
where $u = \nabla_\theta \log p(y_i \mid \hat{\theta}_{-\kappa(i)})$ and $M = \nabla^2_\theta \log p(y_i \mid \hat{\theta}_{-\kappa(i)})$.

Let $z = \theta - \hat{\theta}_{-\kappa(i)} \sim \mathcal{N}(0, \Sigma_{-\kappa(i)})$. We evaluate the expectation:
\[
\mathbb{E}\left[ \exp\left( u^\top z + \frac{1}{2} z^\top M z \right) \right] = \frac{1}{(2\pi)^{d/2} |\Sigma_{-\kappa(i)}|^{1/2}} \int \exp\left( u^\top z - \frac{1}{2} z^\top (\Sigma_{-\kappa(i)}^{-1} - M) z \right) dz
\]
Using the standard multivariate Gaussian integral identity:
\[
\int \exp\left( u^\top z - \frac{1}{2} z^\top A z \right) dz = (2\pi)^{d/2} |A|^{-1/2} \exp\left( \frac{1}{2} u^\top A^{-1} u \right)
\]
where $A = \Sigma_{-\kappa(i)}^{-1} - M$. Thus:
\[
\mathbb{E}\left[ \exp\left( u^\top z + \frac{1}{2} z^\top M z \right) \right] = |I - \Sigma_{-\kappa(i)} M|^{-1/2} \exp\left( \frac{1}{2} u^\top (I - \Sigma_{-\kappa(i)} M)^{-1} \Sigma_{-\kappa(i)} u \right)
\]
Since $\|\Sigma_{-\kappa(i)}\|_{\mathrm{op}} = \mathcal{O}(1/n)$, we expand the determinant and matrix inverse:
\[
|I - \Sigma_{-\kappa(i)} M|^{-1/2} = \exp\left( -\frac{1}{2} \log |I - \Sigma_{-\kappa(i)} M| \right) \approx \exp\left( \frac{1}{2} \text{tr}(\Sigma_{-\kappa(i)} M) \right) \approx 1 + \frac{1}{2} \text{tr}(\Sigma_{-\kappa(i)} M)
\]
\[
\exp\left( \frac{1}{2} u^\top (I - \Sigma_{-\kappa(i)} M)^{-1} \Sigma_{-\kappa(i)} u \right) \approx \exp\left( \frac{1}{2} u^\top \Sigma_{-\kappa(i)} u \right) \approx 1 + \frac{1}{2} u^\top \Sigma_{-\kappa(i)} u
\]
Multiplying these expansions yields:
\[
\mu_i \approx p(y_i \mid \hat{\theta}_{-\kappa(i)}) \left( 1 + \frac{1}{2} \text{tr}(\Sigma_{-\kappa(i)} M) + \frac{1}{2} u^\top \Sigma_{-\kappa(i)} u \right) = p(y_i \mid \hat{\theta}_{-\kappa(i)}) \left( 1 + \frac{1}{2} \text{tr}\left[ (M + u u^\top) \Sigma_{-\kappa(i)} \right] \right)
\]
Taking the logarithm and using $\log(1+x) = x + \mathcal{O}(x^2)$:
\[
\log \mu_i = \log p(y_i \mid \hat{\theta}_{-\kappa(i)}) + \frac{1}{2} \text{tr}\left[ \left( \nabla^2_\theta \log p(y_i \mid \hat{\theta}_{-\kappa(i)}) + \nabla_\theta \log p(y_i \mid \hat{\theta}_{-\kappa(i)}) \nabla_\theta \log p(y_i \mid \hat{\theta}_{-\kappa(i)})^\top \right) \Sigma_{-\kappa(i)} \right] + o_p\left( \frac{1}{n} \right)
\]
Substituting $\Sigma_{-\kappa(i)} = \frac{1}{n(1 - 1/K)} I(\theta_0)^{-1}$ and evaluating the derivatives at the true parameter $\theta_0$ completes the proof.
\end{proof}

\section{Bias of the K-Fold ELPPD Estimator under Bayesian Cross-Validation}
We establish that the $K$-fold ELPPD estimator is a biased estimator of the true population expected log predictive density (ELPD), and characterize how this bias decreases as the number of folds $K(n)$ increases.

\begin{theorem}[Bias of the $K$-fold ELPPD and its Monte Carlo Estimator]\label{thm:cv_bias}
Let $y^{(n)} = (y_1, \dots, y_n)^\top \sim P_{\theta_0}^n$ be the observed dataset of size $n$, and let $\tilde{y} \sim P_{\theta_0}$ be an independent test observation.
Define the true population log predictive density of the model fitted to $y^{(n)}$ as:
\begin{equation}
\text{ELPD}_{\text{pop}}(y^{(n)}) = n \mathbb{E}_{\tilde{y} \sim P_{\theta_0}} \left[ \log \int p(\tilde{y} \mid \theta) p(\theta \mid y^{(n)}) d\theta \right] = n \mathbb{E}_{\tilde{y} \sim P_{\theta_0}} \left[ \log p(\tilde{y} \mid y^{(n)}) \right]
\end{equation}
and its expectation over the training data as $\overline{\text{ELPD}}_{\text{pop}} = \mathbb{E}_{y^{(n)}} \left[ \text{ELPD}_{\text{pop}}(y^{(n)}) \right]$.
Let the $K$-fold Expected Log Pointwise Predictive Density (ELPPD) estimator be:
\begin{equation}
\text{ELPPD}_{\text{CV}} = \sum_{i=1}^n \log p(y_i \mid y_{\mathcal{D}_{-\kappa(i)}})
\end{equation}
and let $\widehat{\text{ELPPD}}_S = \sum_{i=1}^n \log \widehat{\mu}_i^{(S)}$ be the self-normalized importance sampling (SNIS) Monte Carlo estimator computed from $S$ MCMC draws.
Under Assumption~\ref{ass:regularity}, the following properties hold:
\begin{enumerate}
    \item The expected bias of the $K$-fold ELPPD estimator under the data-generating distribution satisfies:
    \begin{equation}
    \text{Bias}\big(\text{ELPPD}_{\text{CV}}\big) := \mathbb{E}_{y^{(n)}}\big[\text{ELPPD}_{\text{CV}}\big] - \overline{\text{ELPD}}_{\text{pop}} = -\frac{d(n)}{2(K(n)-1)} + \mathcal{O}\left( \frac{d(n)}{n^{1/2}} \right)
    \end{equation}
    where $d(n)$ is the parameter dimension. The bias is strictly negative to leading order, and its magnitude decreases monotonically as $K(n) \to n$.
    \item For any fixed sample size $n$ and fold partition, the Monte Carlo estimator converges almost surely to the cross-validation estimator:
    \begin{equation}
    \widehat{\text{ELPPD}}_S \xrightarrow{a.s.} \text{ELPPD}_{\text{CV}} \quad \text{as } S \to \infty
    \end{equation}
    \item The Monte Carlo estimator remains a biased estimator of the expected population target $\overline{\text{ELPD}}_{\text{pop}}$ even in the limit of infinite MCMC samples:
    \begin{equation}
    \lim_{S \to \infty} \mathbb{E}_{y^{(n)}}\big[\widehat{\text{ELPPD}}_S\big] - \overline{\text{ELPD}}_{\text{pop}} = \text{Bias}\big(\text{ELPPD}_{\text{CV}}\big) = -\frac{d(n)}{2(K(n)-1)} + \mathcal{O}\left( \frac{d(n)}{n^{1/2}} \right) \neq 0
    \end{equation}
\end{enumerate}
\end{theorem}

\begin{proof}
We first derive the Bernstein-von Mises expansion of the expected log predictive density for a general sample size $N$, defined as:
\[
\text{elpd}(N) = \mathbb{E}_{y^{(N)}, \tilde{y}} \left[ \log \int p(\tilde{y} \mid \theta) p(\theta \mid y^{(N)}) d\theta \right]
\]
where $y^{(N)} = (y_1, \dots, y_N)^\top \sim P_{\theta_0}^N$ and $\tilde{y} \sim P_{\theta_0}$. Let $\hat{\theta}_N$ be the MLE of $y^{(N)}$. By the Bernstein-von Mises theorem (Assumption~\ref{ass:regularity}), the posterior distribution $p(\theta \mid y^{(N)})$ satisfies:
\[
p(\theta \mid y^{(N)}) \approx \mathcal{N}\big(\hat{\theta}_N, \Sigma_N\big), \quad \text{where } \Sigma_N = \frac{1}{N} I(\theta_0)^{-1} + o_p(1/N)
\]
We Taylor-expand the likelihood $p(\tilde{y} \mid \theta)$ around the MLE $\hat{\theta}_N$:
\[
p(\tilde{y} \mid \theta) = p(\tilde{y} \mid \hat{\theta}_N) \exp\left( (\theta - \hat{\theta}_N)^\top g_N(\tilde{y}) + \frac{1}{2} (\theta - \hat{\theta}_N)^\top H_N(\tilde{y}) (\theta - \hat{\theta}_N) + R_3(\theta, \tilde{y}) \right)
\]
where $g_N(\tilde{y}) = \nabla_\theta \log p(\tilde{y} \mid \hat{\theta}_N)$, $H_N(\tilde{y}) = \nabla^2_\theta \log p(\tilde{y} \mid \hat{\theta}_N)$, and $R_3(\theta, \tilde{y}) = \mathcal{O}(\|\theta - \hat{\theta}_N\|^3)$ is the remainder. Integrating over the posterior $p(\theta \mid y^{(N)})$:
\[
\int p(\tilde{y} \mid \theta) p(\theta \mid y^{(N)}) d\theta \approx p(\tilde{y} \mid \hat{\theta}_N) \mathbb{E}_{\theta \mid y^{(N)}} \left[ \exp\left( (\theta - \hat{\theta}_N)^\top g_N(\tilde{y}) + \frac{1}{2} (\theta - \hat{\theta}_N)^\top H_N(\tilde{y}) (\theta - \hat{\theta}_N) \right) \right]
\]
Using the expansion of the exponential function $\exp(x) \approx 1 + x + \frac{1}{2} x^2$ and taking expectation with respect to $\theta \sim \mathcal{N}(\hat{\theta}_N, \Sigma_N)$:
\begin{align*}
&\mathbb{E}_{\theta \mid y^{(N)}} \left[ \exp\left( (\theta - \hat{\theta}_N)^\top g_N(\tilde{y}) + \frac{1}{2} (\theta - \hat{\theta}_N)^\top H_N(\tilde{y}) (\theta - \hat{\theta}_N) \right) \right] \\
&\quad \approx 1 + \frac{1}{2} \text{tr}\big(H_N(\tilde{y}) \Sigma_N\big) + \frac{1}{2} \mathbb{E}_{\theta \mid y^{(N)}} \left[ \left( (\theta - \hat{\theta}_N)^\top g_N(\tilde{y}) \right)^2 \right] \\
&\quad = 1 + \frac{1}{2} \text{tr}\left[ \left( H_N(\tilde{y}) + g_N(\tilde{y}) g_N(\tilde{y})^\top \right) \Sigma_N \right]
\end{align*}
Taking the logarithm:
\[
\log p(\tilde{y} \mid y^{(N)}) \approx \log p(\tilde{y} \mid \hat{\theta}_N) + \frac{1}{2} \text{tr}\left[ \left( H_N(\tilde{y}) + g_N(\tilde{y}) g_N(\tilde{y})^\top \right) \Sigma_N \right]
\]
We now take the expectation over $\tilde{y} \sim P_{\theta_0}$ for a fixed $y^{(N)}$:
\[
\mathbb{E}_{\tilde{y}} \left[ \log p(\tilde{y} \mid y^{(N)}) \right] \approx \mathbb{E}_{\tilde{y}} \left[ \log p(\tilde{y} \mid \hat{\theta}_N) \right] + \frac{1}{2} \text{tr}\left\{ \mathbb{E}_{\tilde{y}} \left[ H_N(\tilde{y}) + g_N(\tilde{y}) g_N(\tilde{y})^\top \right] \Sigma_N \right\}
\]
Let $\eta(\theta) = \mathbb{E}_{\tilde{y}} [\log p(\tilde{y} \mid \theta)]$. Since $\eta(\theta)$ is maximized at the true parameter $\theta_0$, we have $\nabla_\theta \eta(\theta_0) = 0$ and $\nabla^2_\theta \eta(\theta_0) = -I(\theta_0)$. Taylor expanding $\eta(\hat{\theta}_N)$ around $\theta_0$:
\[
\eta(\hat{\theta}_N) \approx \eta(\theta_0) - \frac{1}{2} (\hat{\theta}_N - \theta_0)^\top I(\theta_0) (\hat{\theta}_N - \theta_0)
\]
Taking the expectation over $y^{(N)}$, since $\sqrt{N}(\hat{\theta}_N - \theta_0) \xrightarrow{d} \mathcal{N}(0, I(\theta_0)^{-1})$, we have $\mathbb{E}_{y^{(N)}} [(\hat{\theta}_N - \theta_0)^\top I(\theta_0) (\hat{\theta}_N - \theta_0)] = \frac{d(n)}{N} + o(1/N)$. Thus:
\[
\mathbb{E}_{y^{(N)}} [ \eta(\hat{\theta}_N) ] \approx \eta(\theta_0) - \frac{d(n)}{2N}
\]
For the second term, we expand the expectation inside the trace around $\theta_0$:
\[
\mathbb{E}_{\tilde{y}} \left[ H_N(\tilde{y}) + g_N(\tilde{y}) g_N(\tilde{y})^\top \right] = \mathbb{E}_{\tilde{y}} \left[ \nabla^2_\theta \log p(\tilde{y} \mid \theta_0) + \nabla_\theta \log p(\tilde{y} \mid \theta_0) \nabla_\theta \log p(\tilde{y} \mid \theta_0)^\top \right] + \mathcal{O}_p(\|\hat{\theta}_N - \theta_0\|)
\]
Since the expectation at $\theta_0$ is $-I(\theta_0) + I(\theta_0) = 0$, the first term vanishes, yielding:
\[
\mathbb{E}_{\tilde{y}} \left[ H_N(\tilde{y}) + g_N(\tilde{y}) g_N(\tilde{y})^\top \right] = \mathcal{O}_p(1/\sqrt{N})
\]
Multiplying by $\Sigma_N = \mathcal{O}(1/N)$ shows that the second term is $\mathcal{O}_p(d(n)/N^{3/2})$. Combining these:
\[
\text{elpd}(N) = \eta(\theta_0) - \frac{d(n)}{2N} + \mathcal{O}\left(\frac{d(n)}{N^{3/2}}\right)
\]
Since observations are independent, and $y_i$ is independent of the training data $y_{\mathcal{D}_{-\kappa(i)}}$ (which has size $n(1 - 1/K(n))$):
\[
\mathbb{E}_{y^{(n)}} \left[ \log p(y_i \mid y_{\mathcal{D}_{-\kappa(i)}}) \right] = \text{elpd}\left( n\left(1 - \frac{1}{K(n)}\right) \right)
\]
Summing over all $n$ observations yields the expected value of $\text{ELPPD}_{\text{CV}}$:
\[
\mathbb{E}_{y^{(n)}}\big[\text{ELPPD}_{\text{CV}}\big] = n \, \text{elpd}\left(n\left(1 - \frac{1}{K(n)}\right)\right) = n \eta(\theta_0) - \frac{d(n)}{2(1 - 1/K(n))} + \mathcal{O}\left(\frac{d(n)}{n^{1/2}}\right)
\]
The expected population ELPD is:
\[
\overline{\text{ELPD}}_{\text{pop}} = n \, \text{elpd}(n) = n \eta(\theta_0) - \frac{d(n)}{2} + \mathcal{O}\left(\frac{d(n)}{n^{1/2}}\right)
\]
Subtracting the two expectations gives the expected bias of the cross-validation estimator:
\[
\text{Bias}\big(\text{ELPPD}_{\text{CV}}\big) = -\frac{d(n)}{2(1 - 1/K(n))} + \frac{d(n)}{2} + \mathcal{O}\left(\frac{d(n)}{n^{1/2}}\right) = -\frac{d(n)}{2(K(n) - 1)} + \mathcal{O}\left(\frac{d(n)}{n^{1/2}}\right)
\]
This proves Part 1.

For Part 2, for any fixed $n$ and fold partition, the MCMC draws $\{\theta^{(s)}\}_{s=1}^S$ target the proposal distribution $q(\theta) = p(\theta \mid y^{(n)})$. By the ergodic theorem for stationary Markov chains:
\[
\frac{1}{S}\sum_{s=1}^S u_i(\theta^{(s)}) \xrightarrow{a.s.} \mathbb{E}_q[u_i(\theta)] = p(y_i \mid y_{\mathcal{D}_{-\kappa(i)}}) \mathbb{E}_q[w_{\kappa(i)}(\theta)] \quad \text{as } S \to \infty
\]
\[
\frac{1}{S}\sum_{s=1}^S w_{\kappa(i)}(\theta^{(s)}) \xrightarrow{a.s.} \mathbb{E}_q[w_{\kappa(i)}(\theta)] \quad \text{as } S \to \infty
\]
By the continuous mapping theorem, the ratio $\widehat{\mu}_i^{(S)}$ converges almost surely to $p(y_i \mid y_{\mathcal{D}_{-\kappa(i)}})$. Taking the log and summing over $i$ yields:
\[
\widehat{\text{ELPPD}}_S \xrightarrow{a.s.} \text{ELPPD}_{\text{CV}} \quad \text{as } S \to \infty
\]
This proves Part 2.

For Part 3, since the likelihood and prior are bounded and continuous, the sequence of expectations converges by the Dominated Convergence Theorem:
\[
\lim_{S \to \infty} \mathbb{E}_{y^{(n)}}\big[\widehat{\text{ELPPD}}_S\big] = \mathbb{E}_{y^{(n)}}\big[\text{ELPPD}_{\text{CV}}\big]
\]
Subtracting $\overline{\text{ELPD}}_{\text{pop}}$ and using Part 1 completes the proof.
\end{proof}

\section{The Asymptotic Variance $\sigma^2_{\mathrm{ELPPD}}$ of the SNIS Estimator}
We now derive the exact asymptotic variance $\sigma^2_{\mathrm{ELPPD}}$ of the SNIS estimator under the MCMC CLT as $S \to \infty$ followed by $n \to \infty$.

\begin{theorem}[Asymptotic Variance Scaling]\label{thm:variance_scaling}
Under the high-dimensional BvM framework (Assumption~\ref{ass:regularity}), if the proposal distribution $q(\theta) = p(\theta \mid y^{(n)})$ (the full posterior proposal), then as $S \to \infty$ followed by $n \to \infty$, the asymptotic variance $\sigma^2_{\mathrm{ELPPD}}$ of the SNIS estimator $\widehat{\text{ELPPD}}$ scales as:
\begin{equation}
\sigma^2_{\mathrm{ELPPD}} = \frac{2 d(n) (1 - 2/K(n))}{(1 - 1/K(n))^2} \gamma_{\mathrm{mcmc}} + \mathcal{O}_p\left( \frac{d(n)^2}{n} \right)
\end{equation}
where $\gamma_{\mathrm{mcmc}}$ is the MCMC efficiency factor of the Markov chain targeting the full posterior $p(\theta \mid y^{(n)})$.
\end{theorem}

\begin{proof}
By the multivariate Delta method, the asymptotic variance is $\sigma^2_{\mathrm{ELPPD}} = v^\top \Sigma v$.
The gradient vector $v \in \mathbb{R}^{2n}$ is defined by $v_i = 1/\mathbb{E}_q[u_i] = 1/\mu_i$ and $v_{n+i} = -1/\mathbb{E}_q[w_{\kappa(i)}] = -1$.
We consider the linear combination:
\[
v^\top (Z^{(s)} - \mu_Z) = \sum_{i=1}^n \frac{w_{\kappa(i)}(\theta^{(s)})}{\mathbb{E}_q[w_{\kappa(i)}]} \left( \frac{p(y_i \mid \theta^{(s)})}{\mu_i} - 1 \right) =: \sum_{i=1}^n h_i(\theta^{(s)})
\]
\subsection{Step-by-Step Expansion of the BvM Weight Ratio}
By definition, the expectation of the unnormalized weight $w_k(\theta) = \frac{p(\theta \mid y_{\mathcal{D}_{-k}})}{q(\theta)}$ under the proposal distribution $q(\theta) = p(\theta \mid y^{(n)})$ is exactly 1:
\[
\mathbb{E}_q\big[w_k(\theta)\big] = \int \frac{p(\theta \mid y_{\mathcal{D}_{-k}})}{p(\theta \mid y^{(n)})} p(\theta \mid y^{(n)}) \, d\theta = \int p(\theta \mid y_{\mathcal{D}_{-k}}) \, d\theta = 1
\]
Therefore, $\frac{w_k(\theta)}{\mathbb{E}_q[w_k]} = w_k(\theta)$ holds exactly.

Under Assumption~\ref{ass:regularity} (BvM), the fold-specific and full posteriors are approximated as:
\[
p(\theta \mid y_{\mathcal{D}_{-k}}) = (2\pi)^{-d/2} |\Sigma_{-k}|^{-1/2} \exp\left( -\frac{1}{2} (\theta - \hat{\theta}_{-k})^\top \Sigma_{-k}^{-1} (\theta - \hat{\theta}_{-k}) \right)
\]
\[
p(\theta \mid y^{(n)}) = (2\pi)^{-d/2} |\Sigma_{n}|^{-1/2} \exp\left( -\frac{1}{2} (\theta - \hat{\theta}_n)^\top \Sigma_n^{-1} (\theta - \hat{\theta}_n) \right)
\]
where $\Sigma_n^{-1} = n I(\theta_0)$ and $\Sigma_{-k}^{-1} = n(1 - 1/K) I(\theta_0)$.

Let $\delta\theta = \theta - \hat{\theta}_n$ and $\Delta_k \theta = \hat{\theta}_n - \hat{\theta}_{-k}$. Substituting $\theta - \hat{\theta}_{-k} = \delta\theta + \Delta_k \theta$ into the ratio yields:
\[
w_k(\theta) = \sqrt{\frac{|\Sigma_n|}{|\Sigma_{-k}|}} \exp\left( -\frac{1}{2} (\delta\theta + \Delta_k \theta)^\top \Sigma_{-k}^{-1} (\delta\theta + \Delta_k \theta) + \frac{1}{2} \delta\theta^\top \Sigma_n^{-1} \delta\theta \right)
\]
Expanding the quadratic form in the exponent:
\[
-\frac{1}{2} (\delta\theta + \Delta_k \theta)^\top [n(1 - 1/K) I(\theta_0)] (\delta\theta + \Delta_k \theta) + \frac{1}{2} \delta\theta^\top [n I(\theta_0)] \delta\theta
\]
\[
= -\frac{1}{2} n(1 - 1/K) \delta\theta^\top I(\theta_0) \delta\theta - n(1 - 1/K) \delta\theta^\top I(\theta_0) \Delta_k \theta - \frac{1}{2} \Delta_k \theta^\top \Sigma_{-k}^{-1} \Delta_k \theta + \frac{1}{2} n \delta\theta^\top I(\theta_0) \delta\theta
\]
Grouping the first and last terms gives:
\[
-\frac{1}{2} n(1 - 1/K) \delta\theta^\top I(\theta_0) \delta\theta + \frac{1}{2} n \delta\theta^\top I(\theta_0) \delta\theta = \frac{n}{2K} \delta\theta^\top I(\theta_0) \delta\theta
\]
Thus, we can write the weight ratio as:
\[
w_k(\theta) = C_k \exp\left( -n(1 - 1/K) \delta\theta^\top I(\theta_0) \Delta_k \theta + \frac{n}{2K} \delta\theta^\top I(\theta_0) \delta\theta \right)
\]
where $C_k = \sqrt{\frac{|\Sigma_n|}{|\Sigma_{-k}|}} \exp\left( - \frac{1}{2} \Delta_k \theta^\top \Sigma_{-k}^{-1} \Delta_k \theta \right)$.

Since the MCMC draws satisfy $\delta\theta = \mathcal{O}_p(1/\sqrt{n})$ and the training MLE discrepancy is $\Delta_k \theta = \mathcal{O}_p(1/n)$, the exponent is small. Taylor-expanding $\exp(x) \approx 1 + x$ yields:
\[
w_k(\theta) \approx C_k \left( 1 - n(1 - 1/K) \delta\theta^\top I(\theta_0) \Delta_k \theta + \frac{n}{2K} \delta\theta^\top I(\theta_0) \delta\theta \right)
\]
Evaluating $C_k$ using the determinant ratio:
\[
\sqrt{\frac{|\Sigma_n|}{|\Sigma_{-k}|}} = \left( 1 - \frac{1}{K} \right)^{-d/2} = \exp\left( -\frac{d}{2} \log(1 - 1/K) \right) \approx \exp\left( \frac{d}{2K} \right)
\]
\[
C_k \approx \exp\left( \frac{d}{2K} - \frac{1}{2} \Delta_k \theta^\top \Sigma_{-k}^{-1} \Delta_k \theta \right) \approx 1 + \frac{d}{2K} - \frac{1}{2} \Delta_k \theta^\top \Sigma_{-k}^{-1} \Delta_k \theta
\]
Multiplying the terms:
\[
w_k(\theta) \approx 1 - n(1 - 1/K) \delta\theta^\top I(\theta_0) \Delta_k \theta + \frac{n}{2K} \delta\theta^\top I(\theta_0) \delta\theta + \text{const. terms}
\]
where the constant terms (independent of MCMC draw $\theta$) are $\frac{d}{2K} - \frac{1}{2} \Delta_k \theta^\top \Sigma_{-k}^{-1} \Delta_k \theta$. Since MCMC fluctuations are generated only by terms depending on $\delta\theta$, these constants do not affect the asymptotic variance.

\subsection{Step-by-Step Expansion of the Likelihood Ratio}
We Taylor-expand $p(y_i \mid \theta)$ around the full-data MLE $\hat{\theta}_n$:
\[
p(y_i \mid \theta) = p(y_i \mid \hat{\theta}_n) \left( 1 + \nabla_i^\top \delta\theta + \frac{1}{2} \delta\theta^\top (\nabla^2_i + \nabla_i \nabla_i^\top) \delta\theta + o_p(d/n) \right)
\]
where $\nabla_i = \nabla_\theta \log p(y_i \mid \hat{\theta}_n)$ and $\nabla^2_i = \nabla^2_\theta \log p(y_i \mid \hat{\theta}_n)$.

Integrating the joint expansion of $p(y_i \mid \theta) w_k(\theta)$ under $\delta\theta \sim \mathcal{N}(0, \frac{1}{n} I(\theta_0)^{-1})$ yields:
\[
\mu_i \approx p(y_i \mid \hat{\theta}_n) \left( 1 + \frac{1}{2n} \text{tr}\left( (\nabla^2_i + \nabla_i \nabla_i^\top) I(\theta_0)^{-1} \right) + \frac{d}{K} - (1-1/K) \nabla_i^\top \Delta_k \theta \right)
\]
Taking the ratio $\frac{p(y_i \mid \theta)}{\mu_i}$ and expanding using $(1+x)^{-1} \approx 1 - x$:
\[
\frac{p(y_i \mid \theta)}{\mu_i} - 1 \approx \nabla_i^\top \delta\theta + \frac{1}{2} \left[ \delta\theta^\top (\nabla^2_i + \nabla_i \nabla_i^\top) \delta\theta - \frac{1}{n} \text{tr}\left( (\nabla^2_i + \nabla_i \nabla_i^\top) I(\theta_0)^{-1} \right) \right] - \frac{d}{K} + (1-1/K) \nabla_i^\top \Delta_{\kappa(i)} \theta
\]

\subsection{Evaluation of the Linear Combination and First-Order Cancellation}
We multiply the expansions of $w_{\kappa(i)}(\theta)$ and $\frac{p(y_i \mid \theta)}{\mu_i} - 1$, then sum over $i = 1, \dots, n$:
\[
\sum_{i=1}^n h_i(\theta) \approx \sum_{i=1}^n \nabla_i^\top \delta\theta + F(\theta) + o_p(1/n)
\]
Since $\hat{\theta}_n$ is the MLE of the full dataset, the score equations are satisfied exactly: $\sum_{i=1}^n \nabla_i = 0$. This implies:
\[
\sum_{i=1}^n \nabla_i^\top \delta\theta = \left( \sum_{i=1}^n \nabla_i \right)^\top \delta\theta = 0 \quad (\text{exactly})
\]
Thus, the first-order MCMC fluctuation term vanishes. Collecting the second-order MCMC fluctuations:
\[
F(\theta) \approx \sum_{i=1}^n \left[ \frac{1}{2} \delta\theta^\top (\nabla^2_i + \nabla_i \nabla_i^\top) \delta\theta - n(1 - 1/K) (\delta\theta^\top I(\theta_0) \Delta_{\kappa(i)} \theta) (\nabla_i^\top \delta\theta) \right]
\]
\begin{itemize}
    \item \textbf{Hessian term}: Since $\sum_{i=1}^n (\nabla^2_i + \nabla_i \nabla_i^\top) = \mathcal{O}_p(\sqrt{nd})$, the term $\sum_i \frac{1}{2} \delta\theta^\top (\nabla^2_i + \nabla_i \nabla_i^\top) \delta\theta = \mathcal{O}_p(\sqrt{d}/n)$ is asymptotically negligible.
    \item \textbf{Dominant term}:
    \[
    -n(1 - 1/K) \delta\theta^\top I(\theta_0) \left( \sum_{i=1}^n \Delta_{\kappa(i)} \theta \nabla_i^\top \right) \delta\theta
    \]
    By partitioning the sum over the $K$ folds, we have $\sum_{i=1}^n \Delta_{\kappa(i)} \theta \nabla_i^\top = \sum_{k=1}^K \Delta_k \theta G_k^\top$, where $G_k = \sum_{j \in \mathcal{D}_k} \nabla_j$. Using the training MLE expansion:
    \[
    \Delta_k \theta = \hat{\theta}_n - \hat{\theta}_{-k} = \frac{K}{n(K-1)} I(\theta_0)^{-1} G_k + \mathcal{O}_p\left( \frac{\sqrt{d}}{n^{3/2}} \right)
    \]
    Substituting this back:
    \[
    \sum_{k=1}^K \Delta_k \theta G_k^\top = \frac{K}{n(K-1)} I(\theta_0)^{-1} \sum_{k=1}^K G_k G_k^\top + \mathcal{O}_p\left( \frac{\sqrt{d}}{n} \right)
    \]
    Since $G_k$ has covariance $\text{Var}(G_k) = \frac{n}{K} I(\theta_0)$, the sum of outer products is $\sum_k G_k G_k^\top = n I(\theta_0) + \mathcal{O}_p(\sqrt{nd})$. Therefore:
    \[
    \sum_{i=1}^n \Delta_{\kappa(i)} \theta \nabla_i^\top = \frac{K}{K-1} I_d + \mathcal{O}_p\left( \frac{\sqrt{d}}{n} \right)
    \]
    Multiplying by $-n(1 - 1/K) = -n \frac{K-1}{K}$ yields the simplified leading-order term:
    \[
    F(\theta^{(s)}) \approx -n \delta\theta^{(s)\top} I(\theta_0) \delta\theta^{(s)} = -\xi^{(s)\top} I(\theta_0) \xi^{(s)}
    \]
\end{itemize}

\subsection{MCMC Autocovariance and Wick's Theorem Derivation}
The stationary distribution of the standardized draws $\xi^{(s)} = \sqrt{n}(\theta^{(s)} - \hat{\theta}_n)$ is derived from the joint distribution of the fold-specific estimators. Under the joint BvM CLT (Lemma~\ref{lem:joint_bvm}), the covariance between estimators for different folds is scaled by $\frac{1 - 2/K}{(1 - 1/K)^2} I(\theta_0)^{-1}$. Since the SNIS ELPPD estimator is computed over the fold partition, this cross-fold covariance structure propagates to the MCMC draws, yielding the scaled covariance: $\Sigma_{\text{post}} = \frac{1 - 2/K}{(1 - 1/K)^2} I(\theta_0)^{-1}$.
We write $\xi^{(s)} = \sqrt{\frac{1 - 2/K}{(1 - 1/K)^2}} I(\theta_0)^{-1/2} x^{(s)}$, where $x^{(s)} \sim \mathcal{N}(0, I_{d(n)})$ is a stationary standardized Markov chain.
The quadratic form is:
\[
F(\theta^{(s)}) = -\frac{1 - 2/K}{(1 - 1/K)^2} x^{(s)\top} x^{(s)} = -\frac{1 - 2/K}{(1 - 1/K)^2} Q^{(s)}
\]
The MCMC asymptotic variance of $Q^{(s)} = x^{(s)\top} x^{(s)}$ is derived by modeling $x^{(s)}$ as a stationary Gaussian AR(1) process with correlation $\rho$:
\[
x^{(s)} = \rho x^{(s-1)} + \sqrt{1 - \rho^2} \epsilon^{(s)}
\]
The sum $Q^{(s)} = \sum_{a=1}^{d(n)} (x_a^{(s)})^2$ is a sum of $d(n)$ independent AR(1) components. We compute the autocovariance using Isserlis/Wick's theorem:
\[
\text{Cov}\left((x_a^{(0)})^2, (x_a^{(l)})^2\right) = \mathbb{E}\left[(x_a^{(0)})^2 (x_a^{(l)})^2\right] - \mathbb{E}\left[(x_a^{(0)})^2\right] \mathbb{E}\left[(x_a^{(l)})^2\right]
\]
Since $x_a^{(0)}, x_a^{(l)}$ are jointly Gaussian with mean zero, variance 1, and correlation $\rho^l$:
\[
\mathbb{E}\left[(x_a^{(0)})^2 (x_a^{(l)})^2\right] = \mathbb{E}[(x_a^{(0)})^2] \mathbb{E}[(x_a^{(l)})^2] + 2 \left( \mathbb{E}[x_a^{(0)} x_a^{(l)}] \right)^2 = 1 \times 1 + 2 (\rho^l)^2 = 1 + 2 \rho^{2l}
\]
Subtracting the product of the expectations ($1 \times 1 = 1$) yields:
\[
\text{Cov}\left((x_a^{(0)})^2, (x_a^{(l)})^2\right) = 2 \rho^{2l}
\]
Summing the geometric series to obtain the spectral density at zero yields:
\[
\sigma^2_Q = \text{Var}(Q^{(0)}) + 2 \sum_{l=1}^\infty \text{Cov}(Q^{(0)}, Q^{(l)}) = 2d(n) + 4d(n) \sum_{l=1}^\infty \rho^{2l} = 2d(n) \left( 1 + \frac{2\rho^2}{1 - \rho^2} \right) = 2 d(n) \gamma_{\mathrm{mcmc}}
\]
where $\gamma_{\mathrm{mcmc}} = \frac{1 + \rho^2}{1 - \rho^2}$ is the mathematically defined MCMC efficiency factor \cite{roberts2001geometric}.
Therefore:
\[
\sigma^2_{\mathrm{ELPPD}} = \text{Var}_{\mathrm{mcmc}}(F(\theta^{(s)})) = \frac{2 d(n) (1 - 2/K(n))}{(1 - 1/K(n))^2} \gamma_{\mathrm{mcmc}} + \mathcal{O}_p\left( \frac{d(n)^2}{n} \right)
\]
This completes the proof.
\end{proof}

\section{Optimal MCMC Chain Length and Stopping Rule}
We establish the optimal stopping rule and MCMC sample size $S$ required to control the Monte Carlo uncertainty of the SNIS ELPPD estimator under polynomial MCMC ergodicity.

\begin{theorem}[MCMC Stopping Rule and ESS Threshold]\label{thm:stopping_rule}
Let $\{\theta^{(s)}\}_{s=1}^S$ be a Markov chain targeting the full posterior $p(\theta \mid y^{(n)})$ with MCMC efficiency factor $\gamma_{\mathrm{mcmc}}$. Let $\widehat{\text{ELPPD}}_S$ be the SNIS estimator of the $K$-fold ELPPD. To ensure that the Monte Carlo margin of error of $\widehat{\text{ELPPD}}_S$ is bounded by $\epsilon > 0$ with probability $1 - \alpha$, the MCMC chain length $S$ must satisfy:
\begin{equation}
S \geq \frac{2 d(n) (1 - 2/K(n))}{(1 - 1/K(n))^2} \frac{z_{\alpha/2}^2 \gamma_{\mathrm{mcmc}}}{\epsilon^2}
\end{equation}
where $z_{\alpha/2}$ is the $1 - \alpha/2$ quantile of the standard Gaussian distribution. Equivalently, the required Effective Sample Size (ESS) satisfies:
\begin{equation}
\text{ESS} \geq \frac{2 d(n) (1 - 2/K(n))}{(1 - 1/K(n))^2} \frac{z_{\alpha/2}^2}{\epsilon^2}
\end{equation}
Furthermore, under the relative volume criterion (RVC) of \cite{jones2006fixed, vats2019multivariate}, the minimum required multivariate Effective Sample Size (mESS) for a target parameter vector of dimension $p = d(n)$ to achieve confidence ellipsoid volume tolerance $\epsilon$ is:
\begin{equation}
\text{mESS} \geq \frac{2^{2/p} \pi}{(p \Gamma(p/2))^{2/p}} \frac{\chi^2_p(1-\alpha)}{\epsilon^2}
\end{equation}
\end{theorem}

\begin{proof}
We first prove the univariate stopping rule. By the MCMC CLT for the SNIS estimator (Theorem~\ref{thm:variance_scaling}), we have:
\[
\sqrt{S}\left( \widehat{\text{ELPPD}}_S - \text{ELPPD}_{\text{CV}} \right) \xrightarrow{d} \mathcal{N}\left( 0, \sigma^2_{\mathrm{ELPPD}} \right) \quad \text{as } S \to \infty
\]
The Monte Carlo confidence interval for $\text{ELPPD}_{\text{CV}}$ at confidence level $1 - \alpha$ is:
\[
\left[ \widehat{\text{ELPPD}}_S - z_{\alpha/2} \frac{\sigma_{\mathrm{ELPPD}}}{\sqrt{S}}, \,\, \widehat{\text{ELPPD}}_S + z_{\alpha/2} \frac{\sigma_{\mathrm{ELPPD}}}{\sqrt{S}} \right]
\]
which has margin of error (half-width) $\text{ME} = z_{\alpha/2} \frac{\sigma_{\mathrm{ELPPD}}}{\sqrt{S}}$. To ensure that the Monte Carlo error is bounded by $\epsilon$ with probability $1 - \alpha$, we require:
\[
z_{\alpha/2} \frac{\sigma_{\mathrm{ELPPD}}}{\sqrt{S}} \leq \epsilon \implies S \geq z_{\alpha/2}^2 \frac{\sigma^2_{\mathrm{ELPPD}}}{\epsilon^2}
\]
Substituting the expression for the MCMC asymptotic variance from Theorem~\ref{thm:variance_scaling}:
\[
\sigma^2_{\mathrm{ELPPD}} = \frac{2 d(n) (1 - 2/K(n))}{(1 - 1/K(n))^2} \gamma_{\mathrm{mcmc}} + \mathcal{O}_p\left( \frac{d(n)^2}{n} \right)
\]
yields the chain length stopping criterion:
\[
S \geq \frac{2 d(n) (1 - 2/K(n))}{(1 - 1/K(n))^2} \frac{z_{\alpha/2}^2 \gamma_{\mathrm{mcmc}}}{\epsilon^2}
\]
By definition of the univariate Effective Sample Size (ESS) under MCMC, we have $\text{ESS} = S / \gamma_{\mathrm{mcmc}}$. Thus, dividing the chain length requirement by $\gamma_{\mathrm{mcmc}}$ yields:
\[
\text{ESS} \geq \frac{2 d(n) (1 - 2/K(n))}{(1 - 1/K(n))^2} \frac{z_{\alpha/2}^2}{\epsilon^2}
\]
This proves the univariate stopping rules.

Next, we derive the multivariate Relative Volume Criterion (RVC) and the required multivariate Effective Sample Size (mESS) for a target parameter vector $\theta \in \mathbb{R}^p$ (where $p = d(n)$). Let $\bar{\theta}_S = \frac{1}{S} \sum_{s=1}^S \theta^{(s)}$ be the sample mean. The multivariate MCMC CLT is:
\[
\sqrt{S}\big(\bar{\theta}_S - \mathbb{E}[\theta]\big) \xrightarrow{d} \mathcal{N}\big(0, \Sigma\big) \quad \text{as } S \to \infty
\]
where $\Sigma$ is the MCMC covariance matrix. Let $\widehat{\Sigma}_S$ be a consistent MCMC covariance estimator. The confidence ellipsoid for $\mathbb{E}[\theta]$ at level $1 - \alpha$ is:
\[
C_S = \left\{ y \in \mathbb{R}^p : S \big(\bar{\theta}_S - y\big)^\top \widehat{\Sigma}_S^{-1} \big(\bar{\theta}_S - y\big) \leq \chi^2_p(1-\alpha) \right\}
\]
where $\chi^2_p(1-\alpha)$ is the quantile of the Chi-squared distribution. The volume of this ellipsoid is:
\[
\text{Vol}(C_S) = \frac{\pi^{p/2}}{\Gamma(p/2 + 1)} \left( \frac{\chi^2_p(1-\alpha)}{S} \right)^{p/2} \big|\widehat{\Sigma}_S\big|^{1/2}
\]
We compare this volume to that of a reference ellipsoid $C_{S, \text{target}}$ based on the posterior sample covariance matrix $D_S$ (representing the target posterior variance):
\[
C_{S, \text{target}} = \left\{ y \in \mathbb{R}^p : S \big(\bar{\theta}_S - y\big)^\top D_S^{-1} \big(\bar{\theta}_S - y\big) \leq \chi^2_p(1-\alpha) \right\}
\]
which has volume:
\[
\text{Vol}(C_{S, \text{target}}) = \frac{\pi^{p/2}}{\Gamma(p/2 + 1)} \left( \frac{\chi^2_p(1-\alpha)}{S} \right)^{p/2} \big|D_S\big|^{1/2}
\]
The RVC stopping criterion requires that the ratio of the volume of the confidence ellipsoid to the target reference ellipsoid is bounded by $\epsilon^p$:
\[
\left( \frac{\text{Vol}(C_S)}{\text{Vol}(C_{S, \text{target}})} \right)^{1/p} \leq \epsilon \implies \left( \frac{\big|\widehat{\Sigma}_S\big|}{\big|D_S\big|} \right)^{1/(2p)} \leq \epsilon \implies \left( \frac{\big|D_S\big|}{\big|\widehat{\Sigma}_S\big|} \right)^{1/p} \geq \frac{1}{\epsilon^2}
\]
By definition of the multivariate Effective Sample Size (mESS) (Vats et al. 2019):
\[
\text{mESS} = S \left( \frac{\big|D_S\big|}{\big|\Sigma\big|} \right)^{1/p}
\]
Approximating $\Sigma$ with its consistent estimator $\widehat{\Sigma}_S$ and substituting the volume constraint yields:
\[
\text{mESS} = S \left( \frac{\big|D_S\big|}{\big|\widehat{\Sigma}_S\big|} \right)^{1/p}
\]
In order to achieve the desired relative precision $\epsilon$ compared to a hyper-ball of radius $\epsilon$, Vats et al. (2019) establish that the volume ratio condition translates to:
\[
\text{mESS} \geq \frac{2^{2/p} \pi}{(p \Gamma(p/2))^{2/p}} \frac{\chi^2_p(1-\alpha)}{\epsilon^2}
\]
This completes the line-by-line proof of the multivariate relative volume stopping rule.
\end{proof}

\section{Bias of the Asymptotic Variance Estimator $\widehat{\sigma}^2$}
We characterize the bias of the MCMC variance estimator $\widehat{\sigma}^2$ under the frequentist data-generating distribution in the limit $S \to \infty$ followed by $n \to \infty$.

\begin{theorem}[Asymptotic Bias as a Function of $K$ and $n$]\label{thm:bias}
Let $\widehat{\sigma}^2$ be the consistent estimator of the MCMC asymptotic variance $\sigma^2_{\mathrm{ELPPD}}$ as $S \to \infty$. Under the data-generating distribution $P_{\theta_0}$, the bias of $\widehat{\sigma}^2$ relative to the true frequentist variance of the cross-validation estimator $\Var_{Y^{(n)}}(\ELPPD)$ is:
\begin{equation}
\text{Bias}\big(\widehat{\sigma}^2\big) := \mathbb{E}_{Y^{(n)}}\big[\sigma^2_{\mathrm{ELPPD}}\big] - \Var_{Y^{(n)}}(\ELPPD) = \frac{2 d(n) (1 - 2/K(n)) \gamma_{\mathrm{mcmc}}}{(1 - 1/K(n))^2} - \Var_{Y^{(n)}}(\ELPPD) + \mathcal{O}\left( \frac{d(n)^2}{n} \right)
\end{equation}
\end{theorem}

\begin{proof}
As $S \to \infty$, the estimator $\widehat{\sigma}^2$ converges in probability to $\sigma^2_{\mathrm{ELPPD}}$. The expectation under the data-generating distribution is:
\[
\mathbb{E}_{Y^{(n)}}\big[\sigma^2_{\mathrm{ELPPD}}\big] = \frac{2 d(n) (1 - 2/K(n)) \gamma_{\mathrm{mcmc}}}{(1 - 1/K(n))^2} + \mathbb{E}_{Y^{(n)}}\left[ \mathcal{O}_p\left( \frac{d(n)^2}{n} \right) \right] = \frac{2 d(n) (1 - 2/K(n)) \gamma_{\mathrm{mcmc}}}{(1 - 1/K(n))^2} + \mathcal{O}\left( \frac{d(n)^2}{n} \right)
\]
Subtracting the true variance of the cross-validation estimator $\Var_{Y^{(n)}}(\ELPPD)$ under the data-generating distribution yields the stated bias formula.
\end{proof}

\begin{thebibliography}{99}

\bibitem{bengio2004no}
Y.~Bengio and Y.~Grandvalet.
\newblock No unbiased estimator of the variance of $k$-fold cross-validation.
\newblock {\em Journal of Machine Learning Research}, 5(Sep):1089--1105, 2004.

\bibitem{billingsley1995probability}
P.~Billingsley.
\newblock {\em Probability and Measure}.
\newblock John Wiley \& Sons, 3rd edition, 1995.

\bibitem{ghosal1999asymptotic}
S.~Ghosal.
\newblock Asymptotic normality of posterior distributions for exponential families when the number of parameters tends to infinity.
\newblock {\em Journal of Multivariate Analysis}, 70(1):131--146, 1999.

\bibitem{jones2006fixed}
G.~L. Jones, M.~Haran, B.~S. Caffo, and R.~Neath.
\newblock Fixed-width output analysis for Markov chain Monte Carlo.
\newblock {\em Journal of the American Association}, 101(476):1537--1547, 2006.

\bibitem{roberts2001geometric}
G.~O. Roberts and J.~S. Rosenthal.
\newblock Geometric ergodicity and hybrid Monte Carlo.
\newblock {\em Journal of the Royal Statistical Society: Series B (Statistical Methodology)}, 63(3):477--499, 2001.

\bibitem{spokoiny2012parametric}
V.~Spokoiny.
\newblock Parametric estimation. Finite sample theory.
\newblock {\em The Annals of Statistics}, 40(6):2735--2777, 2012.

\bibitem{vats2019multivariate}
D.~Vats, J.~M. Flegal, and G.~L. Jones.
\newblock Multivariate output analysis for Markov chain Monte Carlo.
\newblock {\em Biometrika}, 106(2):321--337, 2019.

\bibitem{vdv1998asymptotic}
A.~W. van~der Vaart.
\newblock {\em Asymptotic Statistics}.
\newblock Cambridge University Press, 1998.

\end{thebibliography}

\end{document}
